\documentclass[11pt]{article}
\usepackage[T1]{fontenc}
\usepackage{lmodern}
\usepackage[margin=1in]{geometry}
\usepackage{amsmath,amssymb,amsthm,mathtools}
\usepackage{microtype}
\usepackage[colorlinks=true,linkcolor=blue,citecolor=blue,urlcolor=blue]{hyperref}
\hypersetup{pdftitle={Real spectra for the separated spherical Robin operators},pdfauthor={Kaiwen Jin}}
\newtheorem{theorem}{Theorem}[section]
\newtheorem{lemma}[theorem]{Lemma}
\newtheorem{proposition}[theorem]{Proposition}
\newcommand{\R}{\mathbb R}
\newcommand{\C}{\mathbb C}
\newcommand{\norm}[1]{\lVert#1\rVert}
\DeclareMathOperator{\Dom}{Dom}
\DeclareMathOperator{\Impart}{Im}
\DeclareMathOperator{\Repart}{Re}
\title{Real spectra for the separated spherical Robin operators}
\author{Kaiwen Jin\\\small\texttt{yui.ui.siki@gmail.com}}
\date{October 15, 2026}
\begin{document}
\maketitle
\begin{abstract}
We prove that the separated spherical-strip operators with PT-symmetric
Robin conditions have real spectrum for every half-width $0<a<\pi/2$,
every real boundary parameter and every integer Fourier mode.
The proof establishes strict interlacing between the associated real
Robin and Dirichlet spectra at every index. A positive Laplace
representation of the Fourier coefficients of $\sec^2x$ gives the
comparison for nonzero integer modes, while a first-order factorization
treats the zero mode. Canonical products then give a strict half-plane
sign for the characteristic ratio and a bounded inverse at every
nonreal spectral parameter.
\end{abstract}

\section{Introduction}

Separation of variables in the Laplace--Beltrami problem on a tubular
neighborhood of the sphere's equator leads to regular one-dimensional
operators with complex Robin conditions \cite{S,KS}. The real parts of
the boundary coefficients depend on the strip width. Siegl's question
is whether their spectra remain real at every boundary strength.

More precisely, for $0<a<\pi/2$, $\alpha\in\R$ and $m\in\mathbb Z$,
put
\[
 \beta=\tfrac12\tan a,\qquad
 V_m(x)=\frac{8m^2-3-\cos(2x)}{8\cos^2x}
       =(m^2-\tfrac14)\sec^2x-\tfrac14.
\]
On complex $L^2(-a,a)$ with Lebesgue measure, consider
\begin{equation}\label{eq:operator}
 H_{a,\alpha,m}u=-u''+V_mu
\end{equation}
with domain
\begin{equation}\label{eq:domain}
 \Dom(H_{a,\alpha,m})=
 \left\{u\in H^2(-a,a):
 \begin{aligned}
 u'(a)+(i\alpha+\beta)u(a)&=0,\\
 u'(-a)+(i\alpha-\beta)u(-a)&=0
 \end{aligned}\right\}.
\end{equation}
Both derivatives are coordinate derivatives in $x$.
These are exactly the operators in \cite[equation (2)]{S} and
\cite[\S4.4.2, equations (4.32)--(4.33)]{KS}.
Their PT symmetry is commutation with the antilinear map
$u(x)\mapsto\overline{u(-x)}$.

Krej\v{c}i\v{r}\'{i}k and Siegl proved reality for the comparison obtained
by dropping $V_m$, since $\beta>0$ \cite[Proposition 4.3]{KS}.
They also proved that, for each fixed separated problem, every
eigenvalue with sufficiently large real part is real
\cite[Proposition 4.4]{KS}. Their perturbation argument leaves a finite
part of the spectrum undecided. Reality of that remaining part was
conjectured in \cite[Remark 4.6]{KS} and posed explicitly in \cite{S}.
The following gives an affirmative answer for the exact separated
operators.

\begin{theorem}\label{thm:main}
For every $0<a<\pi/2$, $\alpha\in\R$ and $m\in\mathbb Z$, the operator
\eqref{eq:operator} with domain \eqref{eq:domain} is densely defined
and closed, and
\[
 \sigma(H_{a,\alpha,m})\subset\R.
\]
\end{theorem}

The proof uses the spectra of two real auxiliary problems. Let $r_n$
and $d_n$ be the real Robin and Dirichlet eigenvalues for the same
potential, indexed from zero. We establish
\[
 r_n<d_n<r_{n+1}\qquad(n\geq0).
\]
The second inequality needs more than the ordinary comparison of form
domains. The Fourier coefficients of $\sec^2x$ admit a positive Laplace
representation, which compares the sine and cosine energies on their
full $H^1$ form domains. This proves the required inequality for
$|m|\geq1$. For $m=0$, the real Robin operator factors as $A^*A$,
and its Dirichlet partner is the operator at $m=1$.

The complex boundary determinant is expressed through a meromorphic
ratio whose zeros are $r_n$ and whose poles are $d_n$.
The strict interlacing, together with a positive normalization below
both spectra, forces this ratio to have negative imaginary part in the
upper half-plane. The boundary equation would require it to be real,
so it cannot vanish at a nonreal parameter. Solving the inhomogeneous
ODE then proves the assertion for the entire spectrum.
All estimates are for a fixed separated realization; no parameter-uniform
resolvent bound is used.

Section~\ref{sec:interlace} proves the self-adjoint comparisons, and
Section~\ref{sec:resolvent} derives the characteristic ratio and the
resolvent. We use complex Sobolev spaces and
$\langle u,v\rangle=\int u\overline v$, linear in the first entry.
The norm $\norm{\cdot}_2$ is taken on the interval under discussion.


\section{The auxiliary self-adjoint spectra}\label{sec:interlace}

Fix $0<a<\pi/2$ and put $\beta=\tfrac12\tan a$ throughout.
The real Robin realization of $-\partial_x^2+V_m$ has boundary
conditions $u'(a)+\beta u(a)=0$ and $u'(-a)-\beta u(-a)=0$.
Write its eigenvalues as $r_n(m)$, $n\geq0$, in increasing order.
Write $d_n(m)$ for the eigenvalues of the Dirichlet realization of the
same differential expression. All these operators act in complex $L^2$.
When $m$ is fixed we suppress it in the notation.

\subsection{Forms and parity}

Set $W(x)=\tfrac12\tan x$. Integration by parts gives
\begin{equation}\label{eq:form-positive}
 \int_{-a}^a\bigl(|u'|^2+V_m|u|^2\bigr)\,dx
 +\beta\bigl(|u(a)|^2+|u(-a)|^2\bigr)
 =\norm{u'+Wu}_2^2+m^2\int_{-a}^a\sec^2x\,|u|^2\,dx.
\end{equation}
The Robin form has domain $H^1(-a,a)$; its restriction to
$H^1_0(-a,a)$ is the Dirichlet form. Their form norms, after adding
$\norm{u}_2^2$, are equivalent to the $H^1$ norm, since $W$ and $V_m$
are bounded on the fixed interval. In particular, both forms are closed
and nonnegative.

The associated operators have the asserted $H^2$ domains. Indeed, testing
the form identity against compactly supported functions gives
$-u''+V_mu\in L^2$, hence $u\in H^2$. Integration by parts against all
form-domain functions then gives exactly the free-endpoint Robin
conditions, with the opposite derivative sign at the left endpoint.
Conversely those $H^2$ functions satisfy the form identity. These are
regular real separated Sturm--Liouville problems: the endpoints are
finite and the potential is real and smooth on the closed interval.
They are self-adjoint and have discrete spectra with complete
orthonormal eigenbases; see \cite[Theorems 2.13, 9.6 and 9.10]{T}.
We shall use the form-domain min-max formula
\begin{equation}\label{eq:minmax}
 \lambda_k=\inf_{\substack{F\subset\Dom(q)\\\dim F=k+1}}
       \ \sup_{\substack{u\in F\\\norm{u}_2=1}}q[u],\qquad k\geq0,
\end{equation}
which follows from that eigenbasis and the form version of
\cite[Theorem 4.10]{T}. Formula \eqref{eq:form-positive} also proves
$r_n,d_n\geq0$.

On $(0,a)$ let
\[
 q_R[u]=\int_0^a\bigl(|u'|^2+V_m|u|^2\bigr)\,dx+\beta|u(a)|^2.
\]
Its domain is $H^1(0,a)$ in the even sector, and
$\{u\in H^1(0,a):u(0)=0\}$ in the odd sector. Denote the corresponding
eigenvalues by $R_{e,k}$ and $R_{o,k}$. Denote by $D_{e,k}$ the
Neumann-at-zero, Dirichlet-at-$a$ eigenvalues, and by $D_{o,k}$ the
eigenvalues with Dirichlet conditions at both ends. Write
\[
 q_D[u]=\int_0^a\bigl(|u'|^2+V_m|u|^2\bigr)\,dx
\]
for these Dirichlet forms. Their domains are
$\{u\in H^1(0,a):u(a)=0\}$ in the even sector and $H^1_0(0,a)$ in the
odd sector. They are the restrictions of $q_R$ obtained by imposing
$u(a)=0$. There is no
boundary contribution at zero, where $W(0)=0$, so the half-interval
version of \eqref{eq:form-positive} proves these forms nonnegative too.

\begin{lemma}\label{lem:endpoints}
For every integer $m$,
\[
 R_{e,k}<D_{e,k},\qquad R_{o,k}<D_{o,k},\qquad k\geq0,
\]
and the full-interval spectra have the parity indexing
\begin{equation}\label{eq:parity}
 r_{2k}=R_{e,k},\quad r_{2k+1}=R_{o,k},\qquad
 d_{2k}=D_{e,k},\quad d_{2k+1}=D_{o,k}.
\end{equation}
In particular $r_n<d_n$ for every $n\geq0$.
\end{lemma}

\begin{proof}
A Dirichlet condition at one endpoint is a codimension-one restriction
of the form domain. Formula \eqref{eq:minmax} gives
$\lambda_k\leq\widetilde\lambda_k\leq\lambda_{k+1}$ for the original
and restricted problems. Their spectra cannot meet. For a fixed
condition at the other endpoint, the solution space at a given spectral
parameter is one-dimensional. A common eigenvalue would therefore give
one nonzero solution with both its value and derivative zero at the
changed endpoint, contradicting uniqueness of the scalar initial-value
problem. This also proves simplicity of each separated eigenvalue.
Thus both inequalities are strict. The coefficients extend smoothly
past the endpoints, so uniqueness applies there as well.

Changing the condition at $a$ proves the first two assertions. Changing
it at zero gives
$R_{e,k}<R_{o,k}<R_{e,k+1}$, and likewise
$D_{e,k}<D_{o,k}<D_{e,k+1}$. Even and odd extension, normalized by
$1/\sqrt2$, identify these half-interval problems with the parity
subspaces of the full-interval operators. At zero their $H^2$ domains
match through $u'(0)=0$ and $u(0)=0$, respectively. This proves
\eqref{eq:parity}.
\end{proof}

\subsection{A positive Fourier kernel}

The extra comparisons $D_{e,k}<R_{o,k}$ and $D_{o,k}<R_{e,k+1}$ will
come from a property specific to $\sec^2x$. Put $\omega=\pi/a$.

\begin{lemma}\label{lem:kernel}
For every integer $\ell\geq1$,
\begin{equation}\label{eq:kernel}
 I_\ell:=\int_0^a\sec^2x\cos(\ell\omega x)\,dx
 =(-1)^\ell\int_0^\infty\mu_a(y)e^{-\ell\omega y}\,dy,
\end{equation}
where the positive integrable function $\mu_a$ is
\begin{equation}\label{eq:mu}
 \mu_a(y)=\Impart\sec^2(a+iy)
 =\frac{2\cos a\sin a\cosh y\sinh y}
 {\bigl(\cos^2a\cosh^2y+\sin^2a\sinh^2y\bigr)^2}>0
 \quad(y>0).
\end{equation}
\end{lemma}

\begin{proof}
The function $f(z)=\sec^2z$ is holomorphic in $|\Repart z|<\pi/2$.
Integrate $f(z)e^{i\ell\omega z}$ around the rectangle with vertices
$-a,a,a+iY,-a+iY$. The top integral tends to zero as $Y\to\infty$,
since $f(x+iY)=O(e^{-2Y})$ uniformly for $|x|\leq a$.
On the vertical sides $e^{i\ell\omega(\pm a+iy)}=(-1)^\ell
e^{-\ell\omega y}$ and $f(-a+iy)=\overline{f(a+iy)}$.
The limiting contour identity is
\[
 0=2I_\ell+i(-1)^\ell\int_0^\infty
       \bigl(f(a+iy)-f(-a+iy)\bigr)e^{-\ell\omega y}\,dy
   =2I_\ell-2(-1)^\ell\int_0^\infty\mu_a(y)e^{-\ell\omega y}\,dy.
\]
Direct computation gives \eqref{eq:mu}. It is $O(y)$ at zero and
$O(e^{-2y})$ at infinity, for fixed $a$, so all limiting integrals
converge absolutely.
\end{proof}

\subsection{The nonzero modes}

\begin{proposition}\label{prop:positive-modes}
If $|m|\geq1$, then
\begin{equation}\label{eq:strong-interlace}
 r_n(m)<d_n(m)<r_{n+1}(m)\qquad(n\geq0).
\end{equation}
\end{proposition}

\begin{proof}
Write $\eta=m^2-\tfrac14>0$. First use the normalized half-integer
systems
\[
 c_n(x)=\sqrt{2/a}\cos((n+\tfrac12)\omega x),\qquad
 s_n(x)=\sqrt{2/a}\sin((n+\tfrac12)\omega x),\qquad n\geq0.
\]
Their coefficient map $u_c=\sum b_nc_n\mapsto u_s=\sum b_ns_n$ is
$L^2$-unitary from the form domain $\{u\in H^1:u(a)=0\}$ onto
$\{u\in H^1:u(0)=0\}$ and preserves $\norm{u'}_2^2$.
Here and below these spaces are on $(0,a)$. To see the full-domain
assertion, the four sine and cosine systems used in this proof are
complete by even or odd reflection of the ordinary exponential basis;
see \cite[Theorem 1.3 and Example (1.23)]{T}. For the half-integer
systems, reflect to $(0,2a)$ antisymmetrically for cosines and
symmetrically for sines, retaining the odd frequencies. Integration by
parts gives
\[
 \norm{u'}_2^2=\sum_{n\geq0}((n+\tfrac12)\omega)^2|b_n|^2
\]
on the stated respective domains. Conversely, finiteness of this
weighted sum gives $H^1$ convergence of the partial sums and the stated
endpoint trace. Thus no derivative condition has been added to a form
domain.

The constant term in $V_m=\eta\sec^2x-\tfrac14$ cancels under the map.
Product-to-sum and \eqref{eq:kernel} give
\begin{align}\label{eq:half-transplant}
 q_R[u_s]-q_D[u_c]
 &=\frac{2\eta}{a}\int_0^\infty\mu_a(y)
       \left|\sum_{n\geq0}(-1)^n b_ne^{-(n+1/2)\omega y}\right|^2\,dy
 \notag\\[-2pt]
 &\hspace{1cm}+\frac{2\beta}{a}
       \left|\sum_{n\geq0}(-1)^n b_n\right|^2>0\qquad(u_c\ne0).
\end{align}
For finite sums the potential matrix difference is
$-(2\eta/a)I_{n+j+1}$, in the notation of Lemma~\ref{lem:kernel},
which proves the equality and its sign. For the full form domains,
Cauchy--Schwarz and the weighted coefficient bound give
$\sum|b_n|<\infty$. The endpoint sums are therefore the actual traces;
the Laplace sums converge uniformly after extracting
$e^{-\omega y/2}$. Since $\mu_a$ is integrable, the integral passes to
the limit, while the remaining terms converge in $H^1$.
For strictness, the power series
$\sum(-1)^nb_nt^n$ is nonzero and analytic for $|t|<1$.
It cannot vanish on an interval of $t=e^{-\omega y}\in(0,1)$.
Together with $\mu_a>0$ and $\eta>0$ this makes the integral positive.

Apply \eqref{eq:half-transplant} on the span of the first $k+1$ odd
Robin eigenfunctions, mapped to cosines. The continuous positive form
difference has a positive minimum on its finite-dimensional unit sphere.
Formula \eqref{eq:minmax} then gives
\begin{equation}\label{eq:first-cross}
 D_{e,k}<R_{o,k}\qquad(k\geq0).
\end{equation}
This minimum need not be uniform in $k$.

For the other comparison use the integer systems
$\sqrt{2/a}\cos(n\omega x)$ and $\sqrt{2/a}\sin(n\omega x)$,
$n\geq1$. Their coefficient map sends the mean-zero subspace of
$H^1(0,a)$ onto $H^1_0(0,a)$. It is unitary in $L^2$ and preserves the
derivative norm. Integration by parts proves the weighted coefficient
identity just as above; for the sine system the constant coefficient of
$u'$ is zero because $u(a)-u(0)=0$. Now
\begin{align}\label{eq:integer-transplant}
 q_R[u_c]-q_D[u_s]
 &=\frac{2\eta}{a}\int_0^\infty\mu_a(y)
       \left|\sum_{n\geq1}(-1)^n b_ne^{-n\omega y}\right|^2\,dy
 \notag\\[-2pt]
 &\hspace{1cm}+\frac{2\beta}{a}
       \left|\sum_{n\geq1}(-1)^n b_n\right|^2>0\qquad(u_c\ne0).
\end{align}
Here the potential difference is $(2\eta/a)I_{n+j}$.
The same summability and power-series argument proves the equality and
strictness on the entire domains.

Intersect the first $k+2$ even Robin eigenspace with the mean-zero
subspace. The intersection has dimension at least $k+1$. On any
$(k+1)$-dimensional subspace of it, the difference in
\eqref{eq:integer-transplant} has a positive unit-sphere minimum
$\delta$. The sine image is an admissible Dirichlet subspace, so
\[
 D_{o,k}\leq R_{e,k+1}-\delta<R_{e,k+1}.
\]
Together with \eqref{eq:first-cross} and Lemma~\ref{lem:endpoints},
this gives \eqref{eq:strong-interlace} at every index.
\end{proof}

\subsection{The zero mode}

The coefficient $m^2-\tfrac14$ is negative at $m=0$.
A first-order factorization gives the comparison in this case.

\begin{proposition}\label{prop:zero-mode}
Formula \eqref{eq:strong-interlace} also holds for $m=0$.
\end{proposition}

\begin{proof}
On $L^2(-a,a)$ define
\[
 A=\partial_x+W,\qquad\Dom(A)=H^1(-a,a),\qquad W=\tfrac12\tan x.
\]
This operator is densely defined and closed: convergence of $u_j$ and
$Au_j$ in $L^2$ implies convergence of $u_j'=Au_j-Wu_j$, hence the
weak-derivative identity for their limit. Integration by parts gives
\[
 A^*=-\partial_x+W,\qquad\Dom(A^*)=H^1_0(-a,a).
\]
Indeed, testing an adjoint vector against compactly supported functions
first puts it in $H^1$; the arbitrary two traces of a vector in
$\Dom(A)$ then force its traces to vanish. These conditions also suffice.
Since $W,W'$ are bounded, the product domains are exactly
\begin{align*}
 \Dom(A^*A)&=\{u\in H^2(-a,a):(Au)(-a)=(Au)(a)=0\},\\
 \Dom(AA^*)&=H^2(-a,a)\cap H^1_0(-a,a).
\end{align*}
For example, $u\in H^1$ and $Au\in H^1_0$ imply
$u'=Au-Wu\in H^1$; the converse follows directly. The second product
imposes no extra derivative trace.

The differential expressions are
\[
 A^*A=-\partial_x^2-\tfrac14\sec^2x-\tfrac14,\qquad
 AA^*=-\partial_x^2+\tfrac34\sec^2x-\tfrac14.
\]
Thus $A^*A$ is the real Robin operator at $m=0$, and $AA^*$ is the
Dirichlet operator at $m=1$. The kernel of $A$ is exactly
$\operatorname{span}\{\sqrt{\cos x}\}$, whereas $AA^*\geq\tfrac12$.
For $\lambda>0$ the eigenspaces correspond through $u\mapsto Au$,
with inverse $v\mapsto A^*v/\lambda$. The domain transfer is explicit:
if $A^*Au=\lambda u$, then $Au\in\Dom(A^*)$ and
$A^*(Au)=\lambda u\in\Dom(A)$, so $Au\in\Dom(AA^*)$.
The converse follows from $AA^*v=\lambda v$ and
$\lambda v\in\Dom(A^*)$. Exactly one zero eigenvalue is omitted,
therefore
\[
 r_{n+1}(0)=d_n(1)\qquad(n\geq0).
\]
The two Dirichlet forms at $m=0,1$ have the same domain and differ by
$\int\sec^2x\,|u|^2\,dx\geq\norm{u}_2^2$. Hence
\[
 d_n(0)+1\leq d_n(1)=r_{n+1}(0).
\]
Lemma~\ref{lem:endpoints}, which has no positivity assumption on
$m^2-\tfrac14$, supplies $r_n(0)<d_n(0)$. This proves the assertion.
\end{proof}

\section{The characteristic ratio and the resolvent}\label{sec:resolvent}

Fix $a$ and $m$ again. The result of Section~\ref{sec:interlace} is the
strict ordering $0\leq r_n<d_n<r_{n+1}$ for all $n\geq0$.
We now use this ordering to control the complex Robin boundary equations.
Let $c(x,z)$ and $s(x,z)$ solve $-v''+V_mv=zv$ with initial data
\[
 c(0,z)=1,\quad c'(0,z)=0,\qquad s(0,z)=0,\quad s'(0,z)=1.
\]
Define their right Robin traces and the two products by
\[
 \begin{gathered}
 C(z)=c'(a,z)+\beta c(a,z),\qquad
 S(z)=s'(a,z)+\beta s(a,z),\\
 F_R(z)=C(z)S(z),\qquad F_D(z)=c(a,z)s(a,z).
 \end{gathered}
\]

\subsection{Entire functions and their zeros}

\begin{lemma}\label{lem:products}
The functions $F_R$ and $F_D$ are entire of order at most $1/2$.
Their zeros are exactly the simple zeros $r_n$ and $d_n$, respectively,
with no common zero. The meromorphic ratio
\[
 M(z)=\frac{F_R(z)}{F_D(z)}
\]
satisfies $M(-1)>0$.
\end{lemma}

\begin{proof}
The initial-value problem is based at the interior point zero.
Its solutions and their derivatives are entire in $z$; see
\cite[Theorem 9.1]{T}. To record their growth, choose $k^2=z$ with
$|k|\geq1$ and use the Volterra equation
\[
 c(x,z)=\cos(kx)+\int_0^x\frac{\sin(k(x-t))}{k}V_m(t)c(t,z)\,dt,
 \qquad 0\leq x\leq a.
\]
Multiplication by $e^{-|\Impart k|x}$ and iteration give
\[
 e^{-|\Impart k|x}|c(x,z)|\leq\exp(a\norm{V_m}_\infty/|k|).
\] The analogous equation for $s$
has free term $\sin(kx)/k$. Differentiation in $x$ introduces at most
one additional factor $1+|k|$. Consequently the four boundary functions,
and their products, are bounded by
$K(1+|z|)^K\exp(K\sqrt{|z|})$, for a constant depending on $a,m$.
Near zero the free sine quotient is defined by its entire power series.
This proves the order bound without choosing a global square-root branch.

A zero of $c(a,z)$ or $s(a,z)$ is exactly an eigenvalue of the
corresponding real half-interval Dirichlet problem; a zero of $C$ or $S$
is a half-interval Robin eigenvalue. Thus all zeros are real and the
lists are complete. The identity $cs'-c's=1$ prevents simultaneous
zeros of $c(a,\cdot),s(a,\cdot)$ and of $C,S$.
To verify analytic simplicity, let $f=c$ or $s$ and let a dot denote
$\partial_z$. Differentiating the equation gives, at a real $z$,
\[
 f(a)\dot f'(a)-f'(a)\dot f(a)=-\int_0^a f(x)^2\,dx.
\]
The initial contribution is zero. At a Dirichlet zero this says
$f'(a)\dot f(a)=\int f^2>0$; at a Robin zero it says
$f(a)\dot{(f'+\beta f)}(a)=-\int f^2<0$.
The other factor in each identity is nonzero by uniqueness.
These facts and \eqref{eq:parity} identify the simple zero lists of
$F_R,F_D$. The strict interlacing excludes cancellation between them.

At $z=-1$ neither Dirichlet trace vanishes. For $f=c(\cdot,-1)$ or
$s(\cdot,-1)$, integration of the equation on $(0,a)$ yields
\[
 (f'(a)+\beta f(a))f(a)=q_R[f]+\norm{f}_2^2>0.
\]
The term at zero vanishes in both cases. Nonnegativity of the half form
therefore gives
\[
 \frac{C(-1)}{c(a,-1)}>0,\qquad \frac{S(-1)}{s(a,-1)}>0.
\]
\end{proof}

\subsection{The half-plane sign}

\begin{lemma}\label{lem:half-plane}
For the ratio in Lemma~\ref{lem:products},
\begin{equation}\label{eq:half-plane}
 \Impart M(z)<0\quad(\Impart z>0),\qquad
 \Impart M(z)>0\quad(\Impart z<0).
\end{equation}
\end{lemma}

\begin{proof}
Hadamard factorization for order below one
\cite[Theorem 5.15, p.~57]{M}, applied after shifting the origin to
$z=-1$, gives
\begin{equation}\label{eq:canonical}
 M(z)=M(-1)\,
 \frac{\displaystyle\prod_{n=0}^\infty\frac{r_n-z}{r_n+1}}
      {\displaystyle\prod_{n=0}^\infty\frac{d_n-z}{d_n+1}}.
\end{equation}
The products converge locally uniformly, and there is no nonconstant
exponential factor. The shift is convenient because all $r_n,d_n$ are
nonnegative. In particular, when $m=0$ the zero $r_0=0$ remains as the
factor $-z$; no factor is divided by $r_0$.

Truncate with $K+1$ numerator zeros and $K$ denominator zeros:
\[
 b_K=M(-1)\frac{\prod_{j=0}^{K-1}(d_j+1)}
                      {\prod_{j=0}^{K}(r_j+1)}>0,\qquad
 M_K(z)=-b_K\frac{\prod_{j=0}^{K}(z-r_j)}
                       {\prod_{j=0}^{K-1}(z-d_j)},\qquad K\geq1.
\]
This retains the normalization $M_K(-1)=M(-1)$ and converges locally
uniformly to $M$ away from the real poles. Interlacing gives positive
residues
\[
 \rho_{K,j}=-b_K
 \frac{\prod_{i=0}^{K}(d_j-r_i)}
      {\prod_{0\leq i<K,\ i\ne j}(d_j-d_i)}>0,
 \qquad 0\leq j<K.
\]
Indeed the numerator has $K-j$ negative factors and the denominator
$K-1-j$. Polynomial division now gives a real $a_K$ such that
\[
 M_K(z)=a_K-b_Kz+\sum_{j=0}^{K-1}\frac{\rho_{K,j}}{z-d_j}.
\]
Thus
\[
 \Impart M_K(z)=-\Impart z\left(b_K+
       \sum_{j=0}^{K-1}\frac{\rho_{K,j}}{|z-d_j|^2}\right)<0
 \qquad(\Impart z>0).
\]
The locally uniform limit has $\Impart M\leq0$ in the upper half-plane.
It is nonconstant there: otherwise the entire function
$F_R-c_0F_D$ would vanish identically for some constant $c_0$,
contradicting $F_R(d_0)\ne0$. The identity theorem and open mapping
theorem \cite[Theorems 3.44 and 4.24]{M} therefore make the sign strict.
Finally $M(\overline z)=\overline{M(z)}$ gives the lower-half-plane sign.
\end{proof}

\subsection{The complex boundary equations}

\begin{proof}[Proof of Theorem~\ref{thm:main}]
Fix also $\alpha\in\R$ and take $z\notin\R$.
Evenness of $V_m$ gives $c(-x,z)=c(x,z)$ and $s(-x,z)=-s(x,z)$.
For a homogeneous solution $u=b_ec+b_os$, the two coordinate Robin
conditions are therefore
\begin{equation}\label{eq:matrix}
 \begin{pmatrix}
 C+i\alpha c&S+i\alpha s\\
 -C+i\alpha c&S-i\alpha s
 \end{pmatrix}
 \binom{b_e}{b_o}=0,
\end{equation}
where $c,s$ without arguments are evaluated at $(a,z)$.
The determinant is
\begin{equation}\label{eq:det}
 2\bigl(CS+\alpha^2cs\bigr).
\end{equation}
For nonreal $z$ the factors $c(a,z),s(a,z)$ do not vanish, since their
auxiliary Dirichlet problems are self-adjoint. A zero of
\eqref{eq:det} would consequently require $M(z)=-\alpha^2\in\R$,
contrary to Lemma~\ref{lem:half-plane}.

To obtain the whole-spectrum assertion, let $g\in L^2(-a,a)$ and set
\[
 p(x)=-\int_0^x\bigl[s(x,z)c(t,z)-c(x,z)s(t,z)\bigr]g(t)\,dt.
\]
The integral is oriented when $x<0$. Its kernel vanishes on the diagonal
and its $x$ derivative is one there. Hence
$p''=(V_m-z)p-g$, and boundedness of the fundamental solutions and
their derivatives on the fixed interval gives
\[
 \norm{p}_{H^2(-a,a)}\leq C_{a,m,z}\norm{g}_2.
\]
In particular its endpoint values and derivatives satisfy the same type
of bound. Add $b_ec+b_os$ and solve the matrix in \eqref{eq:matrix}
with right side equal to minus the two Robin traces of $p$.
Its nonzero determinant bounds both constants by a multiple of
$\norm{g}_2$. This constructs a unique $u$ in the exact domain
\eqref{eq:domain} such that
\begin{equation}\label{eq:inverse}
 (H_{a,\alpha,m}-z)u=g,\qquad
 \norm{u}_{H^2(-a,a)}\leq C_{a,\alpha,m,z}\norm{g}_2.
\end{equation}
Uniqueness follows from the homogeneous matrix equation. No uniformity
in the parameters is required.

The domain is dense because it contains $C_c^\infty(-a,a)$.
For completeness it also defines a closed operator. Fix one nonreal
$z$ and write $R_z$ for the inverse just constructed. If $u_j\to u$
and $H_{a,\alpha,m}u_j\to h$ in $L^2$, then
\[
 u_j=R_z(H_{a,\alpha,m}u_j-zu_j)
       \longrightarrow R_z(h-zu)\quad\hbox{in }H^2.
\]
The $L^2$ limit is $u$, and the endpoint traces are continuous in $H^2$.
Thus $u$ belongs to \eqref{eq:domain} and $H_{a,\alpha,m}u=h$.
It follows from \eqref{eq:inverse}, for every $z\notin\R$, that
$z$ lies in the resolvent set. Therefore
$\sigma(H_{a,\alpha,m})\subset\R$ for every $0<a<\pi/2$,
$\alpha\in\R$ and $m\in\mathbb Z$, as asserted.
\end{proof}

\paragraph{AI assistance.}
OpenAI Codex was used in the mathematical exploration, drafting, source
checking and revision of this manuscript. The reader and referee checks
used during preparation were AI-based. Independent mathematical review
is requested.

\begin{thebibliography}{9}
\bibitem{S}
P.~Siegl, \emph{PT-symmetric Laplace--Beltrami operator in the strip on a
sphere}, ESF nonselfadjoint spectral theory problem 2010-10, equation~(2),
\url{https://nsa.fjfi.cvut.cz/problems/01_ESF_2010/2010-10/2010-10_ESF_Siegl.pdf}.

\bibitem{KS}
D.~Krej\v{c}i\v{r}\'{i}k and P.~Siegl,
\emph{PT-symmetric models in curved manifolds},
arXiv:1001.2988v1 (2010),
\url{https://arxiv.org/abs/1001.2988v1}.

\bibitem{T}
G.~Teschl, \emph{Mathematical Methods in Quantum Mechanics: With
Applications to Schr\"odinger Operators}, Graduate Studies in Mathematics
\textbf{99}, American Mathematical Society, 2009.
Author-posted version dated February~12, 2009,
\url{https://www.mat.univie.ac.at/~gerald/ftp/book-schroe/schroe.pdf}.

\bibitem{M}
J.~Murphy, \emph{Introduction to Complex Analysis}, lecture notes for
Math~5351, \url{https://pages.uoregon.edu/jamu/Math5351-lecture-notes.pdf},
accessed October~15, 2026.
\end{thebibliography}
\end{document}
